\documentclass[11pt]{ctexart}
\usepackage[numbers, compress]{natbib}
\usepackage[margin=1in]{geometry}
\usepackage{subcaption}

\usepackage[table]{xcolor}  % colors
\definecolor{Gray}{gray}{0.9}
\definecolor{midgreen}{rgb}{0.1,0.5,0.1}
\definecolor{darkgray}{gray}{0.25}
\definecolor{lightblue}{rgb}{0.25,0.25,0.8}
\definecolor{mydarkblue}{rgb}{0,0.08,0.45}
\usepackage[colorlinks,
linkcolor=lightblue,
anchorcolor=blue,
urlcolor=mydarkblue,
citecolor=lightblue]{hyperref}
% \usepackage[capitalize]{cleverref}
\usepackage{makebox}
\usepackage{amsmath}
\usepackage{amssymb}
\usepackage{mathtools}
\usepackage{amsthm}
\usepackage{setspace}
\usepackage{bm}
\usepackage{bbm}


% page margin preferences
\topmargin -1.5cm
\oddsidemargin -0.04cm
\evensidemargin -0.04cm
\textwidth 16.59cm
\textheight 21.94cm 
\parskip 7.2pt
\parindent 0pt

\input{notations.tex}

%Theorem Enviorments 
\newtheorem{theorem}{Theorem}
\newtheorem{lemma}{Lemma}
\newtheorem{corollary}{Corollary}
\newtheorem{fact}{Fact}
% \newtheorem{fact}[theorem]{Fact}
\newtheorem{assumption}{Assumption}
\newtheorem{proposition}{Proposition}
\newtheorem{definition}{Definition}
\newtheorem{example}{Example}
\newtheorem{remark}{Remark}
\newtheorem{claim}{Claim}
\newtheorem{problem}{Problem}


% Reference to an algorithm, lower-case.
% \def\algref#1{algorithm~\ref{#1}}
% Reference to an algorithm, upper case.
% \def\Algref#1{Algorithm~\ref{#1}}
% \def\twoalgref#1#2{algorithms \ref{#1} and \ref{#2}}
% \def\Twoalgref#1#2{Algorithms \ref{#1} and \ref{#2}}

\def\llama{$\mathtt{Llama}$-$\mathtt{3.1}$-$\mathtt{8B}$-$\mathtt{Instruct}$}
\def\mistral{$\mathtt{Ministral}$-$\mathtt{7B}$-$\mathtt{Instruct}$}
\newcommand\wh[1]{\hstretch{2}{\hat{\hstretch{.5}{#1}}}}
\def\whK{\widehat{\makebox*{$m$}{$ \K $}}}
\def\whV{\widehat{\makebox*{$m$}{$ \V $}}}
\def\bPsi{{\boldsymbol \Psi}}
\def\bpsi{{\boldsymbol \psi}}
\def\bfpn{b_{\text{FPN}}}

\usepackage{xspace}
\newcommand{\TurboQuant}{\textsc{TurboQuant}\xspace}

\newcommand{\email}[1]{\href{mailto:#1}{\color{black} \texttt{#1}}}

\title{TurboQuant：具有近最优失真率的在线向量量化}


\author{
  Amir Zandieh\\
  Google Research\\
  \email{zandieh@google.com} \\
  \and
  Majid Daliri\\
  New York University\\
  \email{daliri.majid@nyu.edu} \\
  \and 
  Majid Hadian\\
  Google DeepMind\\
  \email{majidh@google.com} \\
  \and 
  Vahab Mirrokni \\
  Google Research\\
  \email{mirrokni@google.com} \\
}


\date{}

\usepackage[capitalize]{cleverref}
\usepackage{algorithm}
\usepackage{algorithmic}
% Stabilize references under XeLaTeX + ctex in this translated build.
\let\cref\ref
\let\Cref\ref
% \usepackage[linesnumbered,boxruled,vlined]{algorithm2e}
% \usepackage[vlined]{algorithm2e}
\usepackage{amsthm}

\usepackage{multirow,makecell}
\usepackage{booktabs}

% tikz package
\usepackage{tikz}
\usetikzlibrary{matrix}
\usetikzlibrary{positioning}



\begin{document}

\maketitle
% \footnote{Equal Contribution}

\begin{abstract}
向量量化（Vector Quantization, VQ）起源于 Shannon 源编码理论，目标是在压缩高维欧氏向量时尽量减小其几何结构失真。
本文提出 \\TurboQuant，同时面向均方误差（MSE）与内积失真进行优化，克服现有方法难以达到最优失真率的问题。
我们提出的数据无关算法适用于在线场景，并在任意比特宽度与维度下达到接近最优（仅差小常数因子）的失真率。
其核心做法是：先对输入向量进行随机旋转，使坐标呈现集中化的 Beta 分布；再利用高维下不同坐标近似独立的性质，对每个坐标应用最优标量量化器。
考虑到 MSE 最优量化器在内积估计上会引入偏差，我们提出两阶段方案：先进行 MSE 量化，再对残差施加 1-bit Quantized JL（QJL）变换，从而得到无偏内积量化器。
此外，我们给出了任意向量量化器可达失真率的信息论下界，证明 \\TurboQuant 与该下界仅相差一个较小常数（约 $2.7$）。
实验结果验证了理论分析：在 KV cache 量化中，3.5 bit/通道可实现几乎无质量损失，2.5 bit/通道仅有轻微退化；在近邻搜索中，本方法在召回率上优于现有产品量化方法，并将索引时间降低到几乎为零。
\end{abstract}


\section{引言}
\label{intro}

欧氏空间中的向量量化（Vector Quantization, VQ）是高效处理高维向量的关键技术，广泛应用于从大规模 AI/深度学习模型训练与部署到检索系统向量数据库等多个计算场景。 
其核心目标是在将高维向量压缩（即把浮点坐标量化为低比特整数）时，尽可能降低失真；常见度量包括均方误差（MSE）与内积误差。
在尽量保持这些几何性质后，系统就能以更低延迟、更少计算和通信开销完成内积查询。

这一问题可追溯至 Shannon 在源编码理论中的奠基性工作~\cite{shannon1948mathematical, shannon1959coding}。该理论指出：块编码（即向量量化器）可达到的最小失真由 Shannon 失真率函数决定，而该函数由数据源统计特性和失真度量（如 MSE）共同决定。 
如今，VQ 已成为 AI、深度学习与搜索系统中的基础能力之一。


VQ 的一个关键应用是 AI 模型部署，尤其是大语言模型（LLM）~\cite{achiam2023gpt, dubey2024llama, claude, team2024gemini}。 
LLM 能力高度依赖模型规模和上下文长度~\cite{kaplan2020scaling}，因此服务侧往往面临高内存占用与更高推理时延。 
这类时延主要来自加速器上 HBM 与 SRAM 之间，或分布式集群节点之间的通信瓶颈。 
通过对权重和激活进行压缩/量化，可以有效缓解这些瓶颈并显著降低推理成本。 
而激活与权重之间的内积运算正是深度学习模型的核心，因此量化方案必须在压缩的同时尽量保真地保留内积。 


基于解码器的 Transformer 模型~\cite{vaswani2017attention}也提供了另一个关键应用场景。
在生成过程中，模型需要在 KV cache 中保存历史 token 的 key/value 向量；其规模会随模型大小（层数、头数）与上下文长度同步增长。 
这种增长会在显存占用和计算速度上形成明显瓶颈，尤其在长上下文场景下更为突出。 
因此，如何在不显著牺牲精度的前提下压缩 KV cache 非常关键。
在该任务中，保持嵌入向量的欧氏结构（尤其是内积与距离）对模型性能至关重要。 
VQ 因而成为解决这一挑战的自然框架：它能够压缩高维嵌入，同时保留其关键几何性质。


此外，在高维空间下基于内积或余弦相似度的近邻搜索（NN）~\cite{elasticsearch, guo2020accelerating}是向量数据库的核心能力~\cite{pinecone, gdrant, pgvector}。
这类数据库也是 RAG~\cite{gao2023retrieval, edge2024local}和信息检索~\cite{khattab2020colbert, santhanam2021colbertv2}的重要基础设施。 
VQ（在 NN 文献中常称为 Product Quantization, PQ）在这里同样扮演关键角色。 
它可高效压缩库向量、降低内存压力，并支持对查询向量内积的低延迟近似计算，从而实现快速且准确的近邻检索。

现有 VQ 方法通常存在取舍：要么与加速器向量化不兼容、计算较慢，不适用于 KV cache 量化等实时 AI 场景；要么在比特宽度维度上的失真界不够理想。
本文目标是提出同时克服这些问题的算法。
我们提出 \TurboQuant：其特点是轻量、支持在线应用（对 KV cache 场景尤为关键）、并且高度加速器友好，适配现代 AI 工作负载。 

\TurboQuant 的核心是两阶段设计。 
第一阶段，我们构造在 MSE 意义下具有最优失真率的向量量化器； 
第二阶段，我们对残差施加 1-bit 量化，得到无偏且低失真的内积量化器。 
我们证明仅针对 MSE 优化的量化器并不能保证内积估计无偏，而两阶段方案可以有效弥补这一差距。
在 MSE 量化器中，我们先对 $d$ 维输入向量进行随机旋转。 
利用旋转后各坐标服从 Beta 分布这一关键事实，我们通过求解连续 k-means 问题，为每个坐标构造最优的 Lloyd-Max 量化器~\cite{lloyd1982least, max1960quantizing}。
该方法同时给出最优的 MSE 失真界，并最小化残差的 L2 范数。
为实现无偏且低失真的内积量化，我们再将其与近期提出的 Quantized Johnson-Lindenstrauss（QJL）变换~\cite{qjl}组合，对残差向量各坐标做单比特量化。
最终，本文算法在 MSE 与内积两个指标上都可证明达到近最优失真界，并且在比特宽度依赖上相较已有方法实现了指数级改进。


\subsection{问题定义}

形式化地，我们的目标是设计量化映射 $Q: \RR^d \to \{ 0, 1 \}^B$，将 $d$ 维向量编码为长度为 $B$ 的二进制串。
若令 $B = b \cdot d$（其中 $b \ge 0$），则该量化器的比特宽度为 $b$，即每个实值坐标的平均编码比特数。
同时我们需要逆映射（反量化映射）$Q^{-1}: \{ 0, 1 \}^B \to \RR^d$，用于从量化表示近似重建原向量。
由于 $Q$ 非双射，这一过程天然有损。 
因此，本文的核心目标是最小化失真，重点关注均方误差（MSE）和内积失真。 


我们不对输入向量分布作额外假设，而是采用最坏情形分析。 
并且允许量化器 $Q(\cdot)$ 是随机化的，从而输出具有随机性。
在随机量化器设定下，更合理的度量是对量化随机性取期望后的失真。 
因此，对于任意目标比特宽度 $b$，我们希望对任意（最坏情形）向量 $\x,\y\in\RR^d$ 最小化以下期望失真：
\begin{align}
    \textbf{(MSE)}~~~~~ & D_{\tt mse} := \E_{Q}\left[\left\| \x - Q^{-1}\left( Q(\x) \right) \right\|_2^2 \right] \label{eq:mse}\\
    \textbf{（内积误差）}~~~~~ & D_{\tt prod} := \E_{Q}\left[\left| \langle \y, \x \rangle - \langle \y, Q^{-1}\left( Q(\x) \right) \rangle \right|^2 \right]. \label{eq:prod_error}
\end{align}
上述期望均对量化器 $Q(\cdot)$ 的随机性求取。
此外，对内积量化器我们还要求内积估计无偏，这是许多应用中的关键性质。
更具体地，需要满足：
\begin{align*}
    \textbf{（内积无偏）}~~~~~ &\E_{Q}\left[ \langle \y, Q^{-1}\left( Q(\x) \right) \rangle\right] = \langle \y, \x \rangle.
\end{align*}


我们的目标是设计计算高效的量化器 $Q_{\tt mse}$ 与 $Q_{\tt prod}$，使其在任意给定比特宽度 $b$ 下尽可能达到上述失真的最优界。 
其中，我们特别要求 $Q_{\tt prod}$ 给出无偏内积估计。 
更具体地，给定 $n$ 个实值向量 $x_1, x_2, \ldots x_n \in \RR^d$，我们设计如下基础算子：
\begin{itemize}
    \item \textsc{Quant}：高效量化数据集，得到 $Q(\x_1), Q(\x_2), \ldots, Q(\x_n)$。 
    \item \textsc{DeQuant}：在给定量化数据后，可高效重建任意向量，即计算 $Q^{-1}\left(Q(\x_i)\right)$（$i\in[n]$）。
    % \item \textsc{score}: given the quantized dataset, for any query vector $\y \in \RR^d$ computes the inner product estimators $\langle \y, Q^{-1}\left( Q(\x_i) \right) \rangle $ for every $i \in [n]$ very efficiently.
\end{itemize}


\subsection{相关工作}

\paragraph{VQ 的起源。}
向量量化理论起源于 Shannon 关于可达失真率函数的奠基性工作~\cite{shannon1948mathematical, shannon1959coding}。
1963 年，Zador~\cite{zador1964development} 通过高分辨率分析方法给出了固定速率量化在高码率下的极限操作失真率函数，该结果与 Shannon 失真率函数高度一致。
不过，Zador 的工作并未聚焦可实现算法。
随后 Gersho 的经典论文~\cite{gersho1979asymptotically} 进一步推动了该领域：其普及了高分辨率理论、简化了 Zador 结果、引入了格量化，并提出了影响深远的关键猜想。
尽管理论取得了显著进展，向量量化在早期的工程可用性仍不明确。
最直接的编码方式（暴力最近邻搜索）计算代价过高，限制了 VQ 的实际落地。

\paragraph{在线与离线量化。}
在线（数据无关）量化方法可即时应用，不需要针对数据做额外调参或校准~\cite{dettmers2022gpt3, ashkboos2024quarot, liu2024kivi, shah2024flashattention, han2025polarquant}。
与之相对，离线（数据相关）方法通常需要重度预处理和训练过程，以适配数据分布，因此不适合动态数据场景~\cite{kim2023squeezellm}。
例如，~\cite{frantar2022gptq, lin2024awq, xiao2023smoothquant, chee2023quip} 中的方法会借助二阶（Hessian）信息调节量化映射，这通常需要较高预处理成本，部分方案还需要后处理。


\paragraph{在线 KV Cache 压缩。}

已有多类方法用于压缩 KV cache。
其中一类是结构改造~\cite{shazeer2019fast, ainslie2023gqa, dai2024deepseekmoe}，通过重构 Transformer 来减少需要存储的 key-value 对数量。
另一类方法是剪枝或驱逐冗余/不关键 token~\cite{beltagy2020longformer, zhang2024h2o, liu2024scissorhands, xiao2023efficient, zandieh2024subgen, li2024snapkv, han2025balancekv}。

一种简单且有效的路径是直接量化 KV cache。
围绕该目标已经有多种量化方案~\cite{yue2024wkvquant, yang2024no, dong2024qaq, kang2024gear, zhang2024kv, liu2024kivi, hooper2024kvquant, kim2024lexico, han2025polarquant}。
近期提出的 QJL~\cite{qjl} 基于 sketch 技术给出高效、数据无关的 1-bit 量化，并可对内积查询提供无偏估计。
该方法无需数据调参或自适应，本文在内积失真优化量化器中使用了这一技术。

\paragraph{产品量化（PQ）。}
在欧氏数据集上的近邻（NN）搜索任务中，索引体积往往构成显著内存瓶颈，通常通过量化技术缓解；在 NN 文献中这类方法常被称为产品量化（PQ）。
许多此类算法在建索引阶段依赖 k-means 变体训练量化码本~\cite{jegou2010product, babenko2014additive, ge2013optimized, wang2017survey, guo2020accelerating}。
因此，这些方法由于预处理代价高，通常不适合在线场景。

近期工作~\cite{gao2024practical} 提出了一种基于网格的 PQ 方法，无需预处理。
该方法通过将均匀网格投影到单位球面，再搜索与数据点最近的投影点完成量化。
尽管该方法实践表现优于其理论界（可能因为理论分析较松），但其网格投影与二分搜索过程计算较慢，且由于天然缺乏向量化能力，在 GPU 等加速器上并行效率较差。


% \paragraph{Accelerator-friendly Quantization.}



\subsection{方法概览与贡献}

\paragraph{MSE 优化的 \TurboQuant。}
我们的第一个 VQ 算法面向 \cref{eq:mse} 中定义的 MSE 失真最小化。
为此，我们先对输入向量施加随机旋转，使每个坐标在与具体输入无关的意义下呈现 Beta 分布。
在高维 $d$ 下，由于测度集中与中心极限定理，每个坐标分布会收敛到高斯分布 $\Ncal(1, 1/d)$。 
此外，任意两个不同坐标不仅近似不相关，而且近似独立（这比仅不相关更强）。
这种近似独立性是简化量化器设计的关键。
它使我们能够忽略坐标间相互作用，直接对每个坐标应用最优标量量化，同时仍达到近最优失真。

我们通过 Max-Lloyd 算法求解一维连续 k-means 问题，为 Beta 分布随机变量得到最优标量量化器。 
针对工程上常用的一组位宽，我们离线预计算并存储最优码本，以便 \TurboQuant 在线调用时直接使用。

在 \cref{thrm_mse} 中我们证明：对于任意最坏情形向量 $\x \in \RR^d$（且 $\norm{\x}=1$），$b$-bit 的 MSE 优化 \TurboQuant（记为 $Q_{\tt mse}: \RR^d \to \{ 0, 1 \}^{b \cdot d}$）满足如下失真界：
\begin{itemize}
    \item 对任意 $b \ge 0$，有
    $D_{\tt mse}(Q_{\tt mse}) := \E\left[\left\| \x - Q_{\tt mse}^{-1}\left( Q_{\tt mse}(\x) \right) \right\|_2^2 \right] \le \frac{\sqrt{3} \pi}{2} \cdot \frac{1}{4^{b}}$。
    \item 在低位宽下可给出更精细的数值界。
    具体地，当 $b = 1, 2, 3, 4$ 时，$D_{\tt mse}(Q_{\tt mse}) \approx {\bf 0.36}, {\bf 0.117}, {\bf 0.03}, {\bf 0.009}$。
\end{itemize}

注意，单位范数假设 $\norm{x}_2=1$ 是标准设定且并不强。
对不满足该假设的数据集，可以额外存储每个向量的 $L2$ 范数（浮点精度），并在反量化后据此做缩放恢复。

\paragraph{内积优化的 \TurboQuant。}
我们证明：MSE 最优化量化器在内积估计上通常有偏，因此需要另一套 VQ 方案来实现无偏内积估计。
我们的方案是两阶段算法：先以比目标预算少 1 bit 的方式应用上述 $Q_{\tt mse}$，再对残差施加 QJL~\cite{qjl}。
该方案可证明是无偏的，同时具有近最优的内积误差率。

在 \cref{thrm_prod} 中我们证明：对于任意最坏情形向量 $\x,\y\in\RR^d$（且 $\norm{\x}=1$），$b$-bit 的内积优化 \TurboQuant（记为 $Q_{\tt prod}: \RR^d \to \{ 0, 1 \}^{b \cdot d}$）具有如下性质：
\begin{itemize}
    \item $\E\left[ \left< \y, Q_{\tt prod}^{-1}\left( Q_{\tt prod}(\x) \right) \right> \right] = \langle \y, \x \rangle$
    \item $D_{\tt prod}(Q_{\tt prod}) := \E\left[\left| \langle \y, \x \rangle - \langle \y, Q_{\tt prod}^{-1}\left( Q_{\tt prod}(\x) \right) \rangle \right|^2 \right] \le \frac{\sqrt{3} \pi^2 \cdot \norm{\y}_2^2}{d} \cdot \frac{1}{4^{b}}$ for any $b \ge 0$.
    \item 在低位宽下可给出更精细的数值界。
    具体地，当 $b = 1, 2, 3, 4$ 时，$D_{\tt prod}(Q_{\tt prod}) \approx \frac{\bf 1.57}{d}, \frac{\bf 0.56}{d}, \frac{\bf 0.18}{d}, \frac{\bf 0.047}{d}$。
\end{itemize}


\paragraph{下界。}
在 \cref{thrm_lower_bound} 中，我们结合 Shannon 下界与 Yao 极小极大原理，证明：对任意位宽为 $b$ 的随机量化算法 $Q: \RR^d \to \{ 0, 1 \}^{b \cdot d}$，都存在困难输入 $\x, \y \in \RR^d$（且 $\norm{\x}=1$）使得以下下界成立：
\begin{itemize}
    \item $D_{\tt mse}(Q) := \E\left[\left\| \x - Q^{-1}\left( Q(\x) \right) \right\|_2^2 \right] \ge \frac{1}{4^{b}}$
    \item $D_{\tt prod}(Q) = \E\left[\left| \langle \y, \x \rangle - \langle \y, Q^{-1}\left( Q(\x) \right) \rangle \right|^2 \right] \ge \frac{\norm{\y}_2^2}{d} \cdot \frac{1}{4^{b}}$
\end{itemize}

由这些下界可知，\TurboQuant 的 MSE 失真在可证明意义下最多仅比信息论下界大 $\frac{\sqrt{3}\pi}{2}\approx{\bf 2.7}$ 倍。
值得注意的是，在更低位宽下该常数因子会明显下降。 
例如当 $b=1$ 时，\TurboQuant 的失真仅比最优值高约 ${\bf 1.45}$ 倍；实验结果也验证了其在低比特场景下的高效率。

\paragraph{实验结果。}
在 \cref{sec:exp_valivation} 中，我们从实验上验证了本文的理论失真界，结果表明 \TurboQuant 在多类真实数据集上的观测失真与理论预测高度一致，并且接近所给下界。

此外，在 \cref{sec:niah} 与 \cref{sec:exp_kv_end2end} 中，我们展示了 \TurboQuant 在在线 KV cache 量化中的效果。 
具体而言，在 needle-in-a-haystack 任务上可实现几乎完美的长上下文检索，同时在其他长上下文下游任务中保持高性能，并将 KV cache 压缩到超过 $5\times$。

最后，在 \cref{sec:nn_exp} 中我们将 \TurboQuant 应用于多种高维近邻搜索任务。
\TurboQuant 在召回上持续优于数据相关的 Product Quantization（PQ），同时将索引时间降至接近于零。


\section{预备知识} \label{sec:prelim}
我们使用粗体小写字母（如 $\x,\y$）表示向量，使用粗体大写字母（如 $\M$）表示矩阵。
向量 $\x$ 在坐标区间 $[i,j]$（含端点）上的切片记作 $\x_{i:j}$。对于矩阵 $\M$，其第 $i$ 行记作 $\M_{i,:}$，下文也简写为 $\M_i$。

我们用 $\SS^{d-1}$ 表示 $\RR^d$ 中半径为 $1$ 的超球面。
对随机变量 $x$，其微分熵记为 $h(x)$。
对随机变量 $x,y$，互信息记为 $I(x; y)=h(x)-h(x|y)$。


由于 \TurboQuant 通过随机旋转来缓解最坏输入情形，因此理解超球面随机点的统计性质非常关键。
下面的引理给出了后续分析与算法设计所需的一个核心性质：

\begin{lemma}[超球面随机点的坐标分布]\label{lem_coordinate_distribution}
    对任意正整数 $d$，若 $\x \in \SS^{d-1}$ 在单位超球面上均匀分布，则任意 $j\in[d]$ 的坐标 $\x_j$ 服从如下（缩放/平移形式的）Beta 分布：
    \[
    \x_j \sim f_{X}(x) := \frac{\Gamma(d/2)}{\sqrt{\pi} \cdot \Gamma((d-1)/2)} \left( 1 - x^2 \right)^{(d-3)/2}.
    \]
    在高维下，该 Beta 分布收敛到正态分布，即 $f_{X}(\cdot) \to \Ncal(0, 1/d)$。
\end{lemma}
\begin{proof}
$f_X(x)$ 可表示为：维度 $d-1$、半径 $\sqrt{1-x^2}$ 的球面面积，与维度 $d$ 单位球面体积之比，再乘以缩放因子 $1/\sqrt{1-x^2}$（由勾股关系得到）。
因此，
\[
f_X(x) = \frac{\frac{2 \pi^{(d-1)/2}}{\Gamma((d-1)/2)} \cdot (1-x^2)^{(d-2)/2}}{\frac{2 \pi^{d/2}}{\Gamma(d/2)}} \cdot 1/\sqrt{1-x^2}= \frac{\Gamma(d/2)}{\sqrt{\pi} \cdot \Gamma((d-1)/2)} \left( 1 - x^2 \right)^{(d-3)/2}.
\]
\end{proof}

\subsection{Shannon 失真下界}
\label{sec:slb}
Shannon 下界（SLB）是由 Shannon 有损源编码定理~\cite{shannon1959coding}导出的强有力工具，可为任意有损压缩方案的最优可达失真率提供通用下界。
本文采用的是适用于一般 $d$ 维信源、以均方误差（MSE）为失真度量的 SLB 形式。
\begin{lemma}[SLB]\label{lem_slb}
    设 $\x \in \RR^d$ 为随机向量，其分布为任意概率分布 $p_X$，且微分熵 $h(\x)$ 有限。
    对总比特复杂度 $B\ge 0$，定义 MSE 失真率函数 $D(B)$ 为：
    \[
    D(p_X, B) := \inf \left\{ \E \left[ \norm{\x - \y}_2^2 \right] : I(\x; \y) \le B \right\},
    \]
    其中下确界取遍所有满足 $I(\x;\y)\le B$ 的联合分布（$\y\in\RR^d$ 为重建随机向量），并且 $\E \left[ \norm{\x - \y}_2^2 \right]$ 按该联合分布计算期望 MSE 失真。
    则对任意比特复杂度 $B\ge 0$，有如下 Shannon 下界：
    \[
    D(p_X, B) \ge \frac{d}{2 \pi e} \cdot 2^{(2/d) (h(\x) - B)}.
    \]
\end{lemma}
这是一个经典结果，可由反向高斯测试信道证明（证明见 \cite{cover1999elements}）。
我们的下界推导会用到 SLB 在“单位超球面均匀分布”情形下的一个推论，形式如下：
\begin{lemma}[超球面随机点的 SLB]\label{lem_slb_random_sphere}
设 $\x \in \SS^{d-1}$ 在单位超球面上均匀分布，并按 \cref{lem_slb} 定义总比特复杂度为 $B$ 的 MSE 失真率函数 $D(B)$。
则对任意比特复杂度 $B \ge 0$，有：
\[
D(B) \ge 2^{-2B/d}.
\]
\end{lemma}
\begin{proof}
记 $A_d$ 为超球面 $\SS^{d-1}$ 的面积，则超球面均匀分布的熵满足 $h(\x)=\log_2 A_d$。
代入 \cref{lem_slb} 可得
$D(B) \ge \frac{d}{2 \pi e} \cdot {A_d}^{2/d} \cdot 2^{-2B/d}$。
再利用 Gamma 函数的 Stirling 近似：
$A_d = \frac{2 \pi^{d/2}}{\Gamma(d/2)} \ge \left(\frac{2 \pi e}{d} \right)^{d/2} \cdot \sqrt{\frac{2d}{\pi}} \cdot (1 - O(1/d))$。
将其代回上式即可得到结论。
\end{proof}

% \subsection{Panter and Dite Formula for Optimal Scalar Compressor}

\subsection{QJL：1-bit 内积量化}
如前所述，我们设计了两类 VQ 算法：一类用于最小化 MSE，另一类用于最小化内积误差。
我们证明了：MSE 最优量化器并不必然给出无偏内积估计，尤其在低位宽下偏差会比较明显。
针对内积量化，我们采用两阶段方案。
第一阶段，用比目标预算少 1 bit 的 MSE 最优量化器，以最小化残差的 L2 范数。
第二阶段，对残差应用无偏且最优的单比特量化器。
在该单比特内积量化步骤中，我们采用近期提出的 Quantized Johnson-Lindenstrauss（QJL）算法~\cite{qjl}。QJL 是位宽为 1 的最优内积量化器。下面给出 QJL 的定义及其核心理论保证。
\begin{definition}[QJL]\label{def_qjl}
对任意正整数 $d$，QJL 映射 $Q_{\tt qjl}: \RR^d \to \{ -1, +1 \}^d$ 定义为：
\[
Q_{\tt qjl}(\x) := {\tt sign} \left( \S \cdot \x \right) ~~~~\text{ for any } \x \in \RR^d,
\]
其中 $\S \in \RR^{d \times d}$ 是随机矩阵，其元素独立同分布采样自正态分布 $\Ncal(0, 1)$，且 ${\tt sign}$ 对向量逐坐标作用。
其逆映射/反量化映射 $Q_{\tt qjl}^{-1}: \{ -1, +1\}^d \to \RR^d$ 定义为：
\[
Q_{\tt qjl}^{-1}(\z) := \frac{\sqrt{\pi/2}}{d} \cdot \S^\top \cdot \z ~~~~\text{ for any } \z \in \{ -1, +1\}^d.
\]
\end{definition}

下一条引理重述了~\cite{qjl} 的结论：QJL 既无偏，也具有较小的内积失真：
\begin{lemma}[性能保证：QJL]\label{lem_qjl}
设 $Q_{\tt qjl}$ 与 $Q_{\tt qjl}^{-1}$ 如 \cref{def_qjl} 所定义。
则对任意向量 $\x \in \SS^{d-1}$ 和任意 $\y \in \RR^d$，有：
\begin{itemize}
    \item 无偏性：$\E\left[ \left< \y, Q_{\tt qjl}^{-1}\left( Q_{\tt qjl}(\x) \right) \right> \right] = \langle \y, \x \rangle$。
    \item 方差上界：$\mathtt{Var} \left( \left< \y, Q_{\tt qjl}^{-1}\left( Q_{\tt qjl}(\x) \right) \right> \right) \le \frac{\pi}{2 d} \cdot \norm{\y}_2^2$
\end{itemize}
\end{lemma}
\begin{proof}
无偏性可由 \cite{qjl} 的 Lemma~3.2 直接得到。
下面证明方差上界。设 \cref{def_qjl} 中随机矩阵 $\S$ 的各行为 $\s_1, \s_2, \ldots \s_m$，则有：
    \[
    \left< \y, Q_{\tt qjl}^{-1}\left( Q_{\tt qjl}(\x) \right) \right> = \frac{1}{d} \sum_{i\in[d]} \sqrt{\pi/2} \cdot \s_i^\top \y \cdot {\tt sign}(\s_i^\top \x).
    \]
    由于 $\s_i$ 独立同分布，上式是 $d$ 个 i.i.d. 随机样本的平均，其中
    $z_i := \sqrt{\pi/2} \cdot \s_i^\top \y \cdot {\tt sign}(\s_i^\top \x)$，$i \in [d]$。
    利用 \cite{qjl} 中的 Fact~3.4，可先对单个 $z_i$ 的方差进行上界估计：
    \begin{equation}\label{eq_moments_estimator}
    \mathtt{Var} \left( z_i \right) = \pi/2 \cdot \mathtt{Var} \left( \s_i^\top \y \cdot {\tt sign}(\s_i^\top \x) \right) \le \pi/2 \cdot \E \left[ (\s_i^\top \y)^2 \right] = \pi/2 \cdot \norm{y}_2^2,
    \end{equation}
    其中最后一个等式成立的原因是：$\s_i^\top \y$ 服从均值为 0、方差为 $\|\y\|_2^2$ 的高斯分布。
    于是，$d$ 个 i.i.d. 样本 $z_1,\ldots,z_d$ 的平均值方差满足：
    \[
    \mathtt{Var} \left( \left< \y, Q_{\tt qjl}^{-1}\left( Q_{\tt qjl}(\x) \right) \right> \right) = \frac{1}{d^2} \sum_{i\in[d]} \mathtt{Var} ( z_i ) \le \frac{\pi}{2 d} \cdot \norm{\y}_2^2.
    \]
\end{proof}

\section{\TurboQuant：高性能量化}
我们开发了两种面向不同目标的 VQ 算法。
第一种用于最小化量化重建后的 MSE；
第二种用于无偏内积估计，以消除 MSE 最优量化器天然存在的偏差。
这两种算法将在后续小节中详细介绍。

此外，在 \cref{sec:lower_bound} 中，我们给出了任意向量量化器可达最优失真率的信息论下界。
该分析表明：在所有位宽下，\TurboQuant 与该下界仅相差一个较小常数因子，因而具有近最优性。


\subsection{MSE 最优的 \TurboQuant} \label{sec:mse_turbo_alg}
设 $\x \in \SS^{d-1}$ 是 $d$ 维单位球面上的（最坏情形）向量。
我们的目标是在每个坐标使用 $b$ bit 进行量化，同时最小化 \cref{eq:mse} 定义的重建 MSE。
首先，我们将该向量与随机旋转矩阵 $\BPi \in \RR^{d \times d}$ 相乘进行随机化。
矩阵 $\BPi$ 可由 i.i.d. 高斯随机矩阵做 QR 分解得到。

旋转后的向量 $\BPi \cdot \x$ 在单位球面 $\SS^{d-1}$ 上服从均匀分布。
由 \cref{lem_coordinate_distribution} 可知，其各坐标服从 Beta 分布，且在高维下收敛到高斯分布。
另外，在高维情形下，$\BPi \cdot \x$ 的不同坐标近似独立~\cite{vershynin2018high}，因此可对每个坐标独立应用最优标量量化器。
因此，根据 \cref{lem_coordinate_distribution}，问题可化为：为分布
$f_{X}(x) = \frac{\Gamma(d/2)}{\sqrt{\pi} \cdot \Gamma((d-1)/2)} \left( 1 - x^2 \right)^{(d-3)/2}$（$x \in [-1, 1]$）设计标量量化器。

在已知概率分布条件下，最优标量量化可表述为一维连续 k-means 问题。
具体而言，我们将区间 $[-1,1]$ 划分为 $2^b$ 个簇（桶）。
其最优解满足 Voronoi 划分性质~\cite{lloyd1982least}：当质心按升序排列时，相邻质心中点即为区间分界。
记升序质心为 $c_i$，则标量量化可写成如下 k-means 优化问题：
\begin{equation}\label{eq:continuous_k_means}
    \Ccal(f_X, b) := \min_{-1 \le c_1 \le c_2 \le \ldots \le c_{2^b} \le 1} \sum_{i=1}^{2^b} \int_{\frac{c_{i-1} + c_i}{2}}^{\frac{c_i + c_{i+1}}{2}} |x - c_i|^2 \cdot f_X(x)~dx.
\end{equation}
注意，\cref{eq:continuous_k_means} 中的 $\Ccal(f_X, b)$ 表示位宽为 $b$ 时的最优 MSE 代价；后文将对其上界进行估计，从而得到 \TurboQuant 端到端 MSE 的上界。
\cref{eq:continuous_k_means} 可通过迭代数值法求解到任意精度。
我们对一组工程上常用的位宽 $b$ 进行一次性离线求解，并将结果存储供后续量化直接调用。

例如，在中高维 $d$ 下，当 $f_X(x)$ 已较接近高斯分布时，$b=1,2$ 的最优量化质心分别为
$\left\{ \pm \frac{\sqrt{2/\pi}}{\sqrt{d}} \right\}$ 与 $\left\{ \pm \frac{0.453}{\sqrt{d}}, \pm \frac{1.51}{\sqrt{d}} \right\}$。

因此，量化映射 $Q_{\tt mse}: \RR^d \to \{ 0, 1 \}^{b \cdot d}$ 的过程是：先计算 $\BPi \cdot \x$，再为每个坐标寻找最近质心并存储其索引。
反量化映射 $Q_{\tt mse}^{-1}: \{ 0, 1 \}^{b \cdot d} \to \RR^d$ 则通过索引取回质心，最后乘以 $\BPi^\top$ 旋回原坐标系得到重建向量。
上述流程的伪代码见 \cref{alg_turboquant_mse}。


\begin{algorithm}[t!]
\caption{$\TurboQuant_{\tt mse}$：MSE 优化}\label{alg_turboquant_mse}
\begin{algorithmic}[1]
    \STATE{\bfseries 输入：}维度 $d$ 与比特宽度 $b$

    {\ttfamily\textcolor{blue}{// Global Parameters for Setting up $\TurboQuant_{\tt mse}$}} \\
    \STATE Generate a {\ttfamily\textcolor{blue}{random rotation matrix}} $\BPi \in \RR^{d \times d}$

    \STATE Construct {\ttfamily\textcolor{blue}{codebook}} by finding centroids $c_1, c_2, \ldots c_{2^b} \in [-1, 1]$ that minimize MSE cost in \cref{eq:continuous_k_means} \label{alg:codebook_construction}
   \\\hrulefill
   \STATE{\bf Procedure} $\textsc{Quant}_{\tt mse} ( \x )$
   \STATE $\y \gets \BPi \cdot \x$\label{line_y}
   \STATE ${\tt idx}_j \gets \arg\min_{k \in [2^b]}  \abs{\y_j - c_k}$ for every $j \in [d]$\label{line_idx} \hfill \COMMENT{{\ttfamily\textcolor{blue}{ ${\tt idx}_j$'s are $b$-bit integers}}}
    \STATE {\bf 输出：}${\tt idx}$
   \\\hrulefill
   \STATE{\bf Procedure} $\textsc{DeQuant}_{\tt mse} ({\tt idx})$
   \STATE $\tilde{\y}_j \gets c_{{\tt idx}_j}$ for every $j \in [d]$\label{y_tilde}
   \STATE $\tilde{\x} \gets \BPi^\top \cdot \tilde{\y}$
   \STATE {\bf 输出：}$\tilde{\x}$ 
\end{algorithmic}
\end{algorithm}


现在我们给出 $\TurboQuant_{\tt mse}$ 的主要理论保证。
\begin{theorem}[性能保证：$\TurboQuant_{\tt mse}$]\label{thrm_mse}
对于任意位宽 $b \ge 1$ 与任意向量 $\x \in \SS^{d-1}$，\cref{alg_turboquant_mse} 中的 $\textsc{Quant}_{\tt mse}(\x)$ 会输出索引向量 ${\tt idx} \in [2^b]^d$。
当将该索引向量输入 $\textsc{DeQuant}_{\tt mse}({\tt idx})$ 时，可得到重建向量 $\tilde{\x} \in \RR^d$，并满足以下失真上界：
\begin{itemize}
    \item 定义 $D_{\tt mse} := \E_{\tilde{\x}}[ \norm{\x - \tilde{\x}}_2^2 ]$，则对任意 $b \ge 0$，有 $D_{\tt mse} \le \frac{\sqrt{3} \pi}{2} \cdot \frac{1}{4^{b}}$。
    \item 对小位宽（$b = 1, 2, 3, 4$），可得到更细粒度的失真值：$D_{\tt mse} \approx {\bf 0.36}, {\bf 0.117}, {\bf 0.03}, {\bf 0.009}$。
\end{itemize}
\end{theorem}
\begin{proof}
我们首先证明 $D_{\tt mse} = d \cdot \Ccal(f_X, b)$，其中 $\Ccal(f_X, b)$ 是 \cref{eq:continuous_k_means} 定义的标量量化最优 MSE 代价。
设 $\tilde{\y}$ 如 \cref{alg_turboquant_mse} 第~\ref{y_tilde} 行所定义。
由于 $\BPi$ 是旋转矩阵，有：
$\norm{\x - \tilde{\x}}_2 = \norm{\BPi \cdot \x - \tilde{\y}}_2$.
再记 $\y = \BPi \cdot \x$（见 \cref{alg_turboquant_mse} 第~\ref{line_y} 行），代入 $D_{\tt mse}$ 定义可得：
\begin{align*}
    D_{\tt mse} &= \E [\norm{\y - \tilde{\y}}_2^2] \\
    &= \sum_{j \in [d]} \E\left[ |\y_j - \tilde{\y}_j|^2 \right] \\
    &= \sum_{j \in [d]} \E\left[ |\y_j - c_{{\tt idx}_j}|^2 \right]\\
    &= d \cdot \E\left[ |\y_1 - c_{{\tt idx}_1}|^2 \right] \\
    &= d \cdot \min_{-1 \le c_1 \le c_2 \le \ldots \le c_{2^b} \le 1} \sum_{i=1}^{2^b} \int_{\frac{c_{i-1} + c_i}{2}}^{\frac{c_i + c_{i+1}}{2}} |x - c_i|^2 \cdot f_X(x)~dx \\
    &= d \cdot \Ccal(f_X, b).
\end{align*}
上式第三个等号来自 \cref{alg_turboquant_mse} 第~\ref{y_tilde} 行中 $\tilde{\y}$ 的定义；第四个等号来自 \cref{lem_coordinate_distribution}，即各 $\y_j$ 具有相同分布 $\y_j \sim f_X(\cdot)$。
最后两个等号成立的原因是：在第~\ref{line_idx} 行中，$c_{{\tt idx}_j}$ 被选为每个坐标 $\y_j$ 的最近质心。

接下来我们需要对最优 k-means 代价 $\Ccal(f_X, b)$ 做上界估计。
在中等维度 $d$ 下，$f_X \to \Ncal(0, 1/d)$。
通过数值求解 \cref{eq:continuous_k_means}，当 $b = 1,2,3,4$ 时可得
$\Ccal(f_X, b) \approx \frac{0.36}{d}, \frac{0.117}{d}, \frac{0.03}{d}, \frac{0.009}{d}$。
当位宽更大（$b > 4$）时，可应用 Panter-Dite~\cite{panter1951quantization} 的高分辨率标量量化失真公式，得到：
\[
\Ccal(f_X, b) \le \frac{1}{12} \cdot \left( \int f_X(x)^{1/3}~dx \right)^3 \cdot \frac{1}{4^b} =  \frac{\sqrt{3} \pi }{2 d} \cdot \frac{1}{4^b}.
\]
证明完成。
\end{proof}


\paragraph{码本索引的熵编码。}
\TurboQuant 的效率还可通过对“最近码字索引”进行熵编码进一步提升。 
具体而言，量化后各码字索引出现的概率可写为 $p_\ell := \int_{\frac{c_{\ell-1} + c_\ell}{2}}^{\frac{c_\ell + c_{\ell+1}}{2}} f_X(x)~dx$。
若对索引进行最优编码，则平均比特宽度可接近分布 $\{ p_i \}_{i \in [2^b]}$ 的熵。
这种无损压缩不会改变失真，并且可以“免费”降低比特开销。
其中在 $b=4$ 时下降最明显，$\{ p_i \}_{i \in [2^b]}$ 的熵约为 $3.8$。
基于最优前缀码的精细计算表明，平均比特宽度约可再降低 $5\%$。
不过考虑到收益有限，本文未将其纳入 \TurboQuant 主流程，以保持算法简洁与速度。




\subsection{内积最优的 \TurboQuant} \label{sec:prod_turbo_alg}
对于近邻搜索等关键应用，内积估计的无偏性非常重要。
但 \cref{sec:mse_turbo_alg} 中的 $\TurboQuant_{\tt mse}$ 对查询向量并不保证内积无偏。
为直观说明这一点，考虑位宽为 $b=1$ 的情形。
在该情形下，当 $d$ 足够大时，\cref{eq:continuous_k_means} 的最优码本为 $\left\{ \pm \sqrt{\frac{2}{\pi d}} \right\}$。
这意味着 $\TurboQuant_{\tt mse}$ 的量化映射可写为 $Q_{\tt mse}(\x) = {\tt sign} \left( \BPi \cdot \x \right)$（对任意 $\x \in \RR^d$），对应反量化映射为 $Q_{\tt mse}^{-1}(\z) = \sqrt{\frac{2}{\pi d}} \cdot \BPi^\top \cdot \z$（对任意 $\z \in \{ -1, +1\}^d$）。
因此，当 $d$ 足够大时，由 \cref{lem_qjl} 可得 $\E\left[ \left< \y, Q_{\tt mse}^{-1}\left( Q_{\tt mse}(\x) \right) \right> \right] = \frac{2}{\pi} \cdot \langle \y, \x \rangle$，即存在 $2/\pi$ 的乘性偏差。
如 \cref{sec:exp_valivation} 的实验所示，该偏差会随位宽 $b$ 增加而减弱。


为消除这一偏差，我们提出将 $\TurboQuant_{\tt mse}$ 与 QJL~\cite{qjl} 结合的方案。
具体地，令 $Q_{\tt mse}$ 表示位宽为 $b-1$ 的 $\TurboQuant_{\tt mse}$ 量化映射。
对任意 $\x \in \SS^{d-1}$，定义残差 $\r := \x - Q_{\tt mse}^{-1}\left( Q_{\tt mse}(\x) \right)$。该残差的 L2 范数较小，即期望上有 $\E[\norm{\r}] = \sqrt{\Ccal(f_X, b-1)}$（见 \cref{eq:continuous_k_means}）。
随后对残差应用 QJL 量化映射 $Q_{\tt qjl}$，即可在总位宽为 $b$ 的预算下得到如下无偏内积估计器：
\[
\left< \y, Q_{\tt mse}^{-1}\left( Q_{\tt mse}(\x) \right) \right> + \norm{\r}_2 \cdot \left< \y, Q_{\tt qjl}^{-1}\left( Q_{\tt qjl}(\r) \right) \right>.
\]
更形式化地，量化映射 $Q_{\tt prod}: \SS^{d-1} \to [2^{b-1}]^d \times \{ -1, 1 \}^d \times \RR$ 定义为：
\[
Q_{\tt prod}(\x) = \left[ Q_{\tt mse}(\x), Q_{\tt qjl}\left( \x - Q_{\tt mse}^{-1}\left( Q_{\tt mse}(\x) \right) \right), \norm{\x - Q_{\tt mse}^{-1}\left( Q_{\tt mse}(\x) \right)}_2 \right].
\]
该过程的伪代码见 \cref{alg_turboquant_prod}。

\begin{algorithm}[t!]
\caption{$\TurboQuant_{\tt prod}$：内积优化}\label{alg_turboquant_prod}
\begin{algorithmic}[1]
    \STATE{\bfseries 输入：}维度 $d$ 与比特宽度 $b$

    {\ttfamily\textcolor{blue}{// Global Parameters for Setting up $\TurboQuant_{\tt prod}$}} \\
    \STATE 实例化一个 {\ttfamily\textcolor{blue}{$\TurboQuant_{\tt mse}$}}（位宽为 $b-1$，见 \cref{alg_turboquant_mse}）\label{turbo_mse}
    \STATE Generate a {\ttfamily\textcolor{blue}{random projection matrix}} $\S \in \RR^{d \times d}$ with i.i.d. entries $\S_{i,j} \sim \Ncal(0, 1)$\label{random_s}
    % \STATE Compute the optimal {\ttfamily\textcolor{blue}{k-means cost}} $\Ccal(f_X, b-1)$ as in \cref{eq:continuous_k_means}
   \\\hrulefill
   \STATE{\bf Procedure} $\textsc{Quant}_{\tt prod} ( \x )$
   \STATE ${\tt idx} \gets \textsc{Quant}_{\tt mse}(\x)$
   \STATE $\r \gets \x - \textsc{DeQuant}_{\tt mse}({\tt idx})$ \label{line_residual} \hfill \COMMENT{{\ttfamily\textcolor{blue}{ residual vector}}} 
   \STATE ${\tt qjl} \gets {\tt sign} \left( \S \cdot \r \right)$ \hfill \COMMENT{{\ttfamily\textcolor{blue}{ QJL on residual vector}}}
   \STATE {\bf 输出：}$({\tt idx}, {\tt qjl}, \norm{\r}_2)$
   \\\hrulefill
   \STATE{\bf Procedure} $\textsc{DeQuant}_{{\tt prod}} ( {\tt idx}, {\tt qjl}, \gamma )$
   \STATE $\tilde{\x}_{\tt mse} \gets \textsc{DeQuant}_{\tt mse}({\tt idx})$ \label{line_tilde_x}
   \STATE $\tilde{\x}_{\tt qjl} \gets \frac{\sqrt{\pi/2}}{d} \cdot \gamma \cdot \S^\top \cdot {\tt qjl}$ \label{line_tilde_x_qjl}
   \STATE {\bf 输出：}$\tilde{\x}_{\tt mse} + \tilde{\x}_{\tt qjl}$ \label{output_tilde_x}
\end{algorithmic}
\end{algorithm}

下面给出 $\TurboQuant_{\tt prod}$ 的主要结论。
\begin{theorem}[性能保证：$\TurboQuant_{\tt prod}$] \label{thrm_prod}
对任意位宽 $b \ge 1$ 与任意向量 $\x \in \SS^{d-1}$，\cref{alg_turboquant_prod} 中的过程 $\textsc{Quant}_{\tt prod}(\x)$ 会输出索引向量 ${\tt idx} \in [2^{b-1}]^d$、符号向量 ${\tt qjl} \in \{ -1, 1 \}^d$ 以及非负标量 $\gamma \ge 0$。
将这些量输入到原语 $\textsc{DeQuant}_{\tt prod}({\tt idx}, {\tt qjl}, \gamma)$ 后，可得到重建向量 $\tilde{\x} \in \RR^d$，并且对任意 $\y \in \RR^d$ 满足：
\begin{itemize}
    \item 内积期望无偏：$\E_{\tilde{\x}}\left[ \left< \y, \tilde{\x} \right> \right] = \langle \y, \x \rangle$
    \item 定义内积失真 $D_{\tt prod} := \E_{\tilde{\x}}\left[\left| \langle \y, \x \rangle - \langle \y, \tilde{\x} \rangle \right|^2 \right]$，则对任意 $b \ge 0$ 有
    $D_{\tt prod} \le \frac{\sqrt{3} \pi^2 \cdot \norm{\y}_2^2}{d} \cdot \frac{1}{4^{b}}$。
    \item 对低位宽（$b = 1, 2, 3, 4$），可得到更细粒度的数值界：
    $D_{\tt prod} \approx \frac{\bf 1.57}{d}, \frac{\bf 0.56}{d}, \frac{\bf 0.18}{d}, \frac{\bf 0.047}{d}$。
\end{itemize}
\end{theorem}
\begin{proof}
% Throughout the proof let $Q_{\tt mse}$ denote the $\textsc{TurboQuant}_{\tt mse}$ which is instantiated in line~\ref{turbo_mse} of \cref{alg_turboquant_prod} and $\S$ is the random matrix defined in line~\ref{random_s}.
我们先计算在给定 $\tilde{\x}_{\tt mse}$ 条件下，内积估计 $\langle \y, \tilde{\x} \rangle$ 的条件期望：
\begin{align*}
    \E \left[ \langle \y, \tilde{\x} \rangle | \tilde{\x}_{\tt mse} \right] &= \E_{\tilde{\x}_{\tt qjl}} \left[ \langle \y, \tilde{\x}_{\tt mse} + \tilde{\x}_{\tt qjl} \rangle | \tilde{\x}_{\tt mse} \right]\\
    &= \langle \y, \tilde{\x}_{\tt mse} \rangle + \E_{\tilde{\x}_{\tt qjl}} \left[ \langle \y, \tilde{\x}_{\tt qjl} \rangle | \tilde{\x}_{\tt mse} \right] \\
    &= \langle \y, \tilde{\x}_{\tt mse} \rangle + \langle \y , \r \rangle\\
    &= \langle \y , \x \rangle,
\end{align*}
其中第一步来自算法第~\ref{output_tilde_x} 行对 $\tilde{\x}$ 的定义；
第三步由 \cref{lem_qjl} 得到；最后一步来自残差定义 $\r = \x - \tilde{\x}_{\tt mse}$（见第~\ref{line_residual} 行）。
再由全期望公式可得无条件期望：
$\E_{\tilde{\x}}\left[ \left< \y, \tilde{\x} \right> \right] = \E_{\tilde{\x}_{\tt mse}} \left[\E \left[ \langle \y, \tilde{\x} \rangle | \tilde{\x}_{\tt mse} \right] \right] = \E[\langle \y , \x \rangle] = \langle \y , \x \rangle$，从而证明定理第一条结论。

接着在计算失真时同样对 $\tilde{\x}_{\tt mse}$ 条件化，得到条件失真：
\begin{align*}
    \E\left[\left. \left| \langle \y, \x \rangle - \langle \y, \tilde{\x} \rangle \right|^2 \right| \tilde{\x}_{\tt mse} \right] &= \E_{\tilde{\x}_{\tt qjl}}\left[\left. \left| \langle \y, \x \rangle - \langle \y, \tilde{\x}_{\tt mse} + \tilde{\x}_{\tt qjl} \rangle \right|^2 \right| \tilde{\x}_{\tt mse} \right] \\
    &= \E_{\tilde{\x}_{\tt qjl}}\left[\left. \left| \langle \y, \r \rangle - \langle \y, \tilde{\x}_{\tt qjl} \rangle \right|^2 \right| \tilde{\x}_{\tt mse} \right] \\
    &= \mathtt{Var} \left( \left. \langle \y, \tilde{\x}_{\tt qjl} \rangle \right| \tilde{\x}_{\tt mse} \right) \\
    &\le \frac{\pi}{2d} \cdot \norm{\r}_2^2 \norm{\y}_2^2,
\end{align*}
其中第二步来自 \cref{alg_turboquant_prod} 第~\ref{line_residual} 与第~\ref{line_tilde_x} 行对 $\r$ 和 $\tilde{\x}_{\tt mse}$ 的定义；
第三步使用了 \cref{lem_qjl} 给出的 $\E[\langle \y, \tilde{\x}_{\tt qjl} \rangle] = \langle \y, \r \rangle$；
最后一步使用 \cref{lem_qjl} 的方差上界，并结合第~\ref{line_tilde_x_qjl} 行中 $\tilde{\x}_{\tt qjl}$ 按 $\gamma=\norm{\r}$ 进行重缩放这一事实。

最后由全期望公式并利用 $\r = \x - \tilde{\x}_{\tt mse}$，可得内积失真上界：
\begin{align*}
    D_{\tt prod} &= \E_{\tilde{\x}_{\tt mse}} \left[ \E\left[\left. \left| \langle \y, \x \rangle - \langle \y, \tilde{\x} \rangle \right|^2 \right| \tilde{\x}_{\tt mse} \right] \right]\\
    &\le \frac{\pi}{2 d} \cdot \norm{\y}_2^2 \cdot \E[\norm{\x - \tilde{\x}_{\tt mse}}_2^2]\\
    &= \frac{\pi}{2 d} \cdot \norm{\y}_2^2 \cdot D_{\tt mse}.
\end{align*}
再调用 \cref{thrm_mse} 在位宽 $b-1$ 下的 MSE 界，即得本定理结论。
\end{proof}

\subsection{下界}
\label{sec:lower_bound}

我们通过证明任意压缩算法可达到失真的下界，说明 \TurboQuant 在任意位宽下都实现了“仅差小常数因子”的最优失真率。
下界证明依赖 Yao 极小极大原理。
该原理将“随机算法 + 最坏确定性输入”的下界，转化为“确定性算法 + 随机输入分布”的下界。
随后我们对后者利用 \cref{sec:slb} 中的 Shannon 下界（SLB）导出可达失真率下界。
形式化地，我们证明如下定理。

\begin{theorem}[最优可达压缩失真的下界]
\label{thrm_lower_bound}
对于任意位宽为 $b$ 的随机量化算法 $Q: \SS^{d-1} \to \{ 0, 1 \}^{b \cdot d}$ 及任意重建映射 $Q^{-1}: \{ 0, 1 \}^{b \cdot d} \to \RR^d$，都存在困难输入 $\x \in \SS^{d-1}$，使得：
\[
D_{\tt mse}(Q) := \E\left[\left\| \x - Q^{-1}\left( Q(\x) \right) \right\|_2^2 \right] \ge \frac{1}{4^{b}}.
\]
并且存在 $\y \in \SS^{d-1}$，使得：
\[
D_{\tt prod}(Q) = \E\left[\left| \langle \y, \x \rangle - \langle \y, Q^{-1}\left( Q(\x) \right) \rangle \right|^2 \right] \ge \frac{1}{d} \cdot \frac{1}{4^{b}}
\]
\end{theorem}
\begin{proof}
由 Yao 极小极大原理可知：最坏输入下最优随机压缩算法的期望 MSE（即 $D_{\tt mse}$），等于某个最困难随机分布下最优确定性压缩算法的期望 MSE。
按定义，后者又不小于“输入服从单位超球面均匀分布时”可达到的最优 MSE。

而位宽为 $b$ 的压缩算法在该均匀球面输入下的最优可达 MSE，下界由 \cref{lem_slb_random_sphere} 给出。
因此可得 $D_{\tt mse} \ge \frac{1}{4^b}$。

进一步，由 $D_{\tt mse} \ge \frac{1}{4^b}$ 及 $D_{\tt mse}$ 的定义可得：
\begin{align*}
    D_{\tt mse} &= \sum_{j=1}^d \E \left[\left| \x_j - \left[ Q^{-1}\left( Q(\x) \right) \right]_j \right|^2 \right]\\
    &= \sum_{j=1}^d \E \left[\left| \langle \e_j, \x \rangle - \langle \e_j, Q^{-1}\left( Q(\x) \right) \rangle \right|^2 \right]\\
    &\ge \frac{1}{4^b}.
\end{align*}
由抽屉原理，必存在某个 $j \in [d]$ 使得
$\E \left[\left| \langle \e_j, \x \rangle - \langle \e_j, Q^{-1}\left( Q(\x) \right) \rangle \right|^2 \right] \ge \frac{1}{d} \cdot \frac{1}{4^b}$，
证明完成。
\end{proof}

我们补充说明：对向量量化的\emph{最坏情形}失真，也可通过“球填充（sphere packing）”论证得到可比较下界（且常数通常更大，因为该问题更困难）~\cite{gersho1982structure}。
但对本文分析而言，\cref{thrm_lower_bound} 更稳健也更贴切。
原因在于它给出的是\emph{期望失真}下界而非最坏误差下界，并与 \cref{thrm_mse} 和 \cref{thrm_prod} 的上界形式严格对齐。




\iffalse
\section{KV Cache 量化} \label{sec:polar-quant-kv}

In this section, we describe how PolarQuant can be applied to the KV cache problem and our practical implementation. Formally, given a stream of $(\q_1,\k_1,\v_1),\ldots,(\q_n,\k_n,\v_n)$, where $\q_i,\k_i,\v_i \in \RR^{d}$ are query, key and value embeddings at $i$-th generation step for all $i \in [n]$. Let $\K_{:i}, \V_{:i} \in \RR^{i \times d}$ be matrices defined by stacking $\k_1, \dots, \k_i$ and $\v_1, \dots, \v_j$ in their rows, respectively. The goal is to compute:
\begin{equation}\label{eq:attn-def}
    \text{softmax}\left(\frac{\K_{:i} \cdot \q_i}{\sqrt{d}}\right)^T\cdot \V_{:i}. 
\end{equation} 
% Keeping all of the key-value pairs in the cache is prohibitively expensive, especially for long sequences. Instead, we opt for approximate computation by sampling a few key-value pairs. Specifically, our goal is to construct an algorithm that at every time step $j$ computes an estimator $z_j$ for $\text{Attn}(q_j, K_j, V_j)$ in sublinear in $n$ time and memory. In particular for given precision $\varepsilon$, $z_j$ should satisfy the following error constraint:

For an efficient token generation, the KV cache at $i$-th generation step $(\K_{:i}, \V_{:i})$ are stored in the memory. To reduce the memory space, we invoke PolarQuant (\cref{alg:polar-quant}) on these embeddings. 
Let ${\whK}_{:i}, {\whV}_{:i} \in \RR^{i \times d}$ be their dequantizations using \textsc{DeQuant} procedure in \cref{alg:polar-quant}. Then, we estimate \cref{eq:attn-def} by computing
\begin{align} \label{eq:attn-approx}
\text{softmax}\left({\frac{{\whK}_{:i}  \cdot \q_i}{\sqrt{d}}}\right) \cdot {\whV}_{:i} .
\end{align}
Note that the na\"ive cache requires $d \cdot \bfpn$ memory space to store each $d$-dimensional embedding where $\bfpn$ is the number of bits to represent a single floating-point number. If we quantize $\log_2d$ level angles with $b$ bits each and keep centroids in $\bfpn$ bits, the memory space becomes $\left( \bfpn + (d-1)b\right)$. 
% In other words, when $b=\alpha \cdot \bfpn$ for some $\alpha \in (0,1)$ the compression ratio is $\alpha + \frac1d$. 
For example, \llama{} is represented by $\bfpn=16$ bits and has $d=128$. For $b=3$, we can save the memory space $4.008$ times. In \cref{sec:exp}, the PolarQuant with KV cache marginally degrades the performance of LLMs on various tasks.

\subsection{工程实现} \label{sec:practice}

\fi








\section{实验}\label{sec:exp}
所有实验均在单张 NVIDIA A100 GPU 上完成。
实验部分分为两部分：一部分用于经验性验证理论结果；另一部分用于评估方法在下游任务中的表现，重点包括 KV cache 量化与近邻向量搜索。

\subsection{经验验证}
\label{sec:exp_valivation}

\begin{figure}[th]
\begin{center}
    \begin{subfigure}{\textwidth}
        \caption{$\TurboQuant_\text{\tt prod}$}
        \centering
        \includegraphics[width=\textwidth]{experiments/analysis/compute_pairwise_inner_product_error_unbiased_1536.pdf}
    \end{subfigure}
    \begin{subfigure}{\textwidth}
        \caption{$\TurboQuant_\text{\tt mse}$}
        \centering
        \includegraphics[width=\textwidth]{experiments/analysis/compute_pairwise_inner_product_error_1536.pdf}
    \end{subfigure}
    \caption{$\TurboQuant_\text{\tt prod}$ 与 $\TurboQuant_\text{\tt mse}$ 在内积估计任务上的误差分布。}
    \label{fig:inner_distortion}
\end{center}
\end{figure}

本节验证前文建立的理论结果。我们使用 DBpedia Entities 数据集开展实验，并采用 OpenAI3 嵌入将其编码到 1536 维空间。
实验中，我们从数据集中随机采样 100,000 条数据作为训练集（主数据集）；另外抽取 1,000 条互不重复的数据作为查询集。

我们评估两种量化方法：\(\TurboQuant_\text{\tt prod}\) 与 \(\TurboQuant_\text{\tt mse}\)。
其中 \(\TurboQuant_\text{\tt mse}\) 面向量化向量与原始向量之间 MSE 的优化；而 \(\TurboQuant_\text{\tt prod}\) 面向内积估计并保证无偏。

\begin{figure}[h]
\begin{center}
    \begin{subfigure}{\textwidth}
        \caption{$\TurboQuant_{\tt prod}$}
        \centering
        \includegraphics[width=\textwidth]{experiments/analysis/ip_distortion_groups_unbiased.pdf}
    \end{subfigure}
    \begin{subfigure}{\textwidth}
        \caption{$\TurboQuant_{\tt mse}$}
        \centering
        \includegraphics[width=\textwidth]{experiments/analysis/ip_distortion_groups.pdf}
    \end{subfigure}
    \caption{在 $b=2$ 时，$\TurboQuant_{\tt prod}$ 的内积误差方差基本保持常数；而 $\TurboQuant_{\tt mse}$ 的方差会随平均内积增大而上升。}
    \label{fig:inner_prod_2bit}
\end{center}
\end{figure}

我们在训练集上执行量化，并将两种方法都用于内积估计任务，分析不同位宽下的内积失真。
如 \cref{fig:inner_distortion} 所示，随着位宽增加，两种方法的方差都下降。
但在内积估计场景下，\(\TurboQuant_\text{\tt mse}\) 会引入偏差；该偏差会随位宽增大而减小，并最终趋近于 0。

\cref{fig:inner_distortion} 的实验结果进一步确认：\(\TurboQuant_\text{\tt prod}\) 在所有位宽下都保持内积无偏，而 \(\TurboQuant_\text{\tt mse}\) 会随着位宽增加逐步改善。

如 \cref{fig:inner_prod_2bit} 所示，在 2 位量化下，$\TurboQuant\text{\tt prod}$ 的方差与原向量内积大小基本无关；而 $\TurboQuant\text{\tt mse}$ 的偏差与平均内积相关，平均内积越大，偏差越大。

\begin{figure}[h]
    \begin{center}
    \begin{tabular}{cc}
        \begin{subfigure}{0.4\textwidth}
            \caption{\textbf{内积误差}}
            \centering
            \includegraphics[width=\textwidth]{experiments/analysis/validation_upper_lower_bounds_innerproduct_unbiased.pdf}
        \end{subfigure}&
        \begin{subfigure}{0.4\textwidth}
            \caption{\textbf{MSE}}
            \centering
            \includegraphics[width=\textwidth]{experiments/analysis/validation_upper_lower_bounds_mse_unbiased.pdf}
    \end{subfigure}
    \label{fig:lower_upper}
    \end{tabular}
    \caption{不同位宽比例下，内积误差与 MSE 相对理论上界/下界的对比。}
    \end{center}
\end{figure}

除直方图外，我们还在 \cref{fig:lower_upper} 中给出了不同位宽比例下原向量与量化向量之间的平均内积误差与 MSE，并与理论分析得到的上界和下界同图对比。
结果显示实验与理论预测一致：在较低位宽比例时，$\TurboQuant\text{\tt prod}$ 在内积估计上更优；而随着位数增加，$\TurboQuant\text{\tt mse}$ 的偏差持续下降，并最终在内积估计上取得更好表现。


\subsection{大海捞针（Needle-In-A-Haystack）} \label{sec:niah}

\begin{figure}[h]
    \centering
    \begin{tabular}{ccc}
        \begin{minipage}{0.33\textwidth}
            \centering
            \textbf{SnapKV} \\ Score: 0.858 \\ % Caption above
            \includegraphics[width=\textwidth]{experiments/needle/needle_snapkv.pdf}
        \end{minipage} &
        \begin{minipage}{0.33\textwidth}
            \centering
            \textbf{PyramidKV} \\ Score: 0.895 \\ % Caption above
            \includegraphics[width=\textwidth]{experiments/needle/needle_pyramidkv.pdf}
        \end{minipage} &
        \begin{minipage}{0.33\textwidth}
            \centering
            \textbf{KIVI} \\ Score: 0.981 \\ % Caption above
            \includegraphics[width=\textwidth]{experiments/needle/needle_kivi.pdf}
        \end{minipage}\vspace{0.5cm} \\
        \begin{minipage}{0.33\textwidth}
            \centering
            \textbf{PolarQuant} \\ Score: 0.995 \\ % Caption above
            \includegraphics[width=\textwidth]{experiments/needle/needle_adapkvqrndrot.pdf}
        \end{minipage} &
        \begin{minipage}{0.33\textwidth}
            \centering
            \textbf{Full-Precision} \\ Score: 0.997 \\ % Caption above
            \includegraphics[width=\textwidth]{experiments/needle/needle_exact.pdf}
        \end{minipage} &
        \begin{minipage}{0.33\textwidth}
            \centering
            \textbf{\TurboQuant} \\ Score: 0.997 \\ % Caption above
            \includegraphics[width=\textwidth]{experiments/needle/needle_rqjl.pdf}
        \end{minipage} \\
    \end{tabular}
    \caption{\texttt{Llama-3.1-8B-Instruct} 在“Needle-In-A-Haystack”测试上的表现。该测试要求模型从长上下文中找回隐藏句子。尽管部分方法召回明显下降，\TurboQuant 在超过 \textbf{$4 \times$} 量化压缩下仍与未压缩基线达到同等性能。}
    \label{fig:needle_comparison}
\end{figure}

% Next we evaluate our method for the ``Needle-In-A-Haystack" test~\cite{niah}. It asks the model to retrieve the information in a given sentence where the sentence (the ``needle'') is placed in an arbitrary location of a long document (the ``haystack''). We follow the same setting from~\citet{fu2024data} and use the \llama{} to run the test. 
“Needle-In-A-Haystack Test”~\cite{niah} 用于评估模型在长文档中检索特定信息的能力。测试做法是在大段文本（haystack）中随机位置插入一条唯一句子（needle），并检测模型能否准确找回。

按照 \citet{fu2024data} 的实验设置，我们使用 \llama{} 进行评测。为分析不同输入长度下的表现，文档长度从 \emph{4k} 到 \emph{104k tokens} 变化。主要评价指标为 \emph{recall score}，用于衡量模型找回隐藏句子的准确程度。

对比方法包括多种 SOTA 内存高效方案：PolarQuant~\cite{han2025polarquant}、SnapKV~\cite{li2024snapkv}、PyramidKV~\cite{cai2024pyramidkv} 与 KIVI~\cite{liu2024kivi}。所有方法都在 0.25 的内存压缩率下测试，即仅使用完整 KV cache 的 25\%。

\cref{fig:needle_comparison} 表明：具有理论保证的量化方法（如 PolarQuant 与 \TurboQuant）优于 token 级压缩方法（SnapKV、PyramidKV）以及缺乏严格理论保证的标量量化方法（如 KIVI）。特别地，\TurboQuant 在 $4\times$ 压缩下仍与全精度模型表现一致，体现出其在长上下文处理中的稳健性。


% We vary the input sequence lengths from 4K to 104K. 
% The evaluation is based on the recall score by comparing the hidden sentence. We compare PolarQuant to SnapKV~\cite{li2024snapkv}, PyramidKV~\cite{cai2024pyramidkv} and KIVI~\cite{liu2024kivi}, where we use their implementations from~\cite{jiang2024minference}. 
% We set hyperparameters to get a memory compression ratio of $0.25$. 
% All methods are set to a compression ratio of $0.25$, i.e., required memory is $\times 0.25$ the full KV cache.  
% In \cref{fig:needle_comparison}, we observe that quantization methods (e.g., KIVI, PolarQuant) outperform token-level compression methods (e.g., SnapKV, PyramidKV). PolarQuant shows better scores than KIVI. Additionally, PolarQuant shows a marginally better score than PolarQuant-R. 

\subsection{LongBench 端到端生成}
\label{sec:exp_kv_end2end}

我们在 LongBench 数据集~\cite{bai2023longbench} 上评测多种 KV cache 压缩算法。该数据集覆盖单文档/多文档问答、摘要、少样本学习、合成任务与代码补全等多类长文本场景。为在不同上下文长度上获得更均衡评估，我们使用长度分布更均匀的子集 \textbf{LongBench-E}，从而更公平地比较各方法在不同上下文规模下的性能。

我们在 \llama{} 与 \mistral{} 两个模型上，将 \TurboQuant 与 \cref{sec:niah} 中的主流基线进行比较。不同于 KIVI、PolarQuant 等通常不量化生成阶段 token 的做法，我们的方法在流式生成过程中也执行量化。

如 \cref{tab:performance_comparison} 所示，在 \llama{} 与 \mistral{} 上，我们的方法均优于其他方案并取得更高平均分。我们在生成过程中评测了 \textbf{2.5-bit} 与 \textbf{3.5-bit} 两种量化配置。
这两种非整数比特精度来自“离群通道/普通通道”分组策略：对两组分别应用独立的 \TurboQuant，并给离群通道分配更高位宽。该策略与已有工作~\cite{zandieh2024qjl, su2025rotatekvaccuraterobust2bit} 一致。
例如在 2.5-bit 配置下，32 个离群通道使用 3 bit，剩余 96 个通道使用 2 bit，有效位宽为 \( (32 \times 3 + 96 \times 2) / 128 = 2.5 \)。3.5-bit 情况采用不同的离群比例，从而获得更高有效位宽。
尽管使用位宽低于若干对比方法，\TurboQuant 仍能保持接近未量化模型的性能，并实现至少 $4.5\times$ 的向量压缩。

% We run various KV cache compression algorithms for LongBench-E datasets~\cite{bai2023longbench}, which encompasses diverse long-text scenarios including single/multi-document question-answering, summarization , few-shot learning, synthetic tasks, and code completion (Code).
% We evaluate \TurboQuant against the top baseline methods using in \cref{sec:niah} on \llama{} and \mistral{}. Despite the fact that for both method KIVI and PolarQuant we didn't quantize the generated tokens, for our method we quantized even generated tokens in the stream. 


% As reported in \cref{tab:performance_comparison}, our methods achieve superior performance compared to other methods, i.e., the average performance scores are higher by a large margin. We run our method for both 2.5 bits and 3.5 bits quantization during the generation and despite the fact the number of bits we're using is smaller compare to other methods, we got unquantized result even though we quantized vectors by factor of at least $\times 4.5$. 

\setlength{\tabcolsep}{12pt}
\def\arraystretch{0.8}%
\begin{table}[t]
\centering
\scalebox{0.9}{
\small
\setlength{\tabcolsep}{4pt}
\begin{tabular}{@{}lcccccccc@{}}
        \toprule 
        \textbf{Method} & \textbf{KV Size} & \textbf{SingleQA} & \textbf{MultiQA} & \textbf{Summarization} & \textbf{Few shot} & \textbf{Synthetic} & \textbf{Code} &\textbf{Average} \\ 
        \midrule
        \multicolumn{9}{c}{\llama{}} \\
        Full Cache    & $16$ & $45.29$ & $45.16$    & $26.55$    & $68.38$   & $59.54$    & $46.28$ & ${50.06}$ \\
        \midrule
        \midrule
        KIVI & $3$ & $43.38$ & $37.99$ & $27.16$    & $68.38$   & $59.50$    & $44.68$ & $48.50$ \\ 
        \midrule
        KIVI & $5$ & $45.04$ & $45.70$ & $26.47$    & $68.57$   & $59.55$    & $46.41$ & $50.16$ \\ 
        \midrule
        \midrule
        PolarQuant & $3.9$ & $45.18$ & $44.48$ & $26.23$    & $68.25$   & $60.07$    & $45.24$ & $49.78$ \\ 
        \midrule
        \midrule
        \TurboQuant (ours) & $2.5$ & $44.16$ & $44.96$ & $24.80$    & $68.01$   & $59.65$    & $45.76$ & $49.44$ \\ 
        \midrule
        \TurboQuant (ours) & $3.5$ & $45.01$ & $45.31$ & $26.00$    & $68.63$   & $59.95$    & $46.17$ & $50.06$ \\ 
        \midrule
        \multicolumn{9}{c}{\mistral{}} \\
        \midrule
        Full Cache    & $16$ & $47.53$ & $49.06$    & $26.09$    & $66.83$   & $53.50$    & $47.90$ & ${49.89}$ \\
        \midrule
        \midrule
        \TurboQuant (ours) & $2.5$ & $48.38$ & $49.22$ & $24.91$    & $66.69$   & $53.17$  & $46.83$ & $49.62$ \\ 
    \end{tabular}
}
\caption{LongBench-V1~\cite{bai2023longbench} 上不同 KV cache 压缩方法在 \llama{} 上的结果。}
\label{tab:performance_comparison}
\end{table}

\subsection{近邻搜索实验}
\label{sec:nn_exp}
本节展示我们的方法在近邻搜索场景中的性能优势。实验使用 DBpedia~\cite{thakur2021beir} Entities 数据集，并基于 OpenAI3 嵌入构建 1536 维\footnote{https://huggingface.co/datasets/Qdrant/dbpedia-entities-openai3-text-embedding-3-large-1536-1M}与 3072 维\footnote{https://huggingface.co/datasets/Qdrant/dbpedia-entities-openai3-text-embedding-3-large-3072-1M}向量空间。此外，我们还在较低维场景下使用标准 GloVe~\cite{pennington2014glove} 嵌入进行评估。
实验设置上，我们随机采样 100,000 个数据点作为训练集（同时用于主要评测）。对于未显式提供查询集的数据集，额外抽取 1,000 个互异样本作为查询点；对 GloVe 数据集则采用其已有的 10,000 条查询集。

我们将 \TurboQuant 与两种基线量化方法比较：Product Quantization（PQ）和 RabitQ~\cite{gao2024practical}。为保证公平性，三种方法都在同一训练集上完成量化，并以 top-k 召回率（1@k）作为评价指标，即衡量真实 top-1 内积结果被包含在近似 top-k 结果中的频率。

\begin{table}[ht]
    \centering
    \begin{tabular}{lccc}
        \toprule
        Approach & d=200 & d=1536 & d=3072 \\
        \midrule
        Product Quantization & 37.04 & 239.75 & 494.42 \\
        RabitQ & 597.25 & 2267.59 & 3957.19 \\
        \TurboQuant & 0.0007 & 0.0013 & 0.0021 \\
        \bottomrule
    \end{tabular}
    \caption{在 4-bit 量化设置下，不同方法在不同维度上的量化时间（秒）。}
    \label{tab:quantizing_time}
\end{table}

\textbf{Product Quantization (PQ)} 依赖 k-means 构造码本，并且需要额外存储码本。
随着 bit 数增加，码本规模呈指数增长，从而带来额外存储开销。在实验中，我们仔细调参以匹配各方法的比特分配。
面向高速查询的高效实现通常使用 AVX2 寄存器内查表（LUT）。
具体来说，LUT16（$l=16$）是速度较快的配置，但我们观察到该配置会带来明显精度下降。
为了兼顾速度与精度，我们采用了 LUT256（256 个码字）的 PQ 版本。
在 2-bit 量化下每次查表覆盖 4 个坐标，在 4-bit 量化下每次查表覆盖 2 个坐标。
此外，由于我们在同一数据集上同时训练和评测，PQ 在该设定下天然占有一定优势。

\textbf{RabitQ。}
与 PQ 不同，RabitQ 缺乏完整向量化实现，无法有效利用 GPU 加速，因此在 CPU 上速度明显更慢。
此外，该方法还有一些额外计算开销，这些开销未被显式计入 bit ratio 比较。也就是说，RabitQ 声称的位宽比在实际系统中往往会因实现开销而高于报告值。

尽管基线方法在设定上有一定优势，\TurboQuant 在全部实验中的召回率仍持续优于 PQ 与 RabitQ。这表明我们的方法在高维量化检索任务中兼具稳健性与效率，是具有竞争力的替代方案。

\begin{figure}[th]
    \centering
    \begin{tabular}{ccc}
        \begin{subfigure}{0.3\textwidth}
            \centering
            \caption{GloVe - d=200}
            \centering
            \includegraphics[width=\textwidth]{experiments/nearest_neighbor/recall-200.pdf}
        \end{subfigure}&
        \begin{subfigure}{0.3\textwidth}
        \centering
            \caption{OpenAI3（$d=1536$）}
            \centering
            \includegraphics[width=\textwidth]{experiments/nearest_neighbor/recall-1536.pdf}
        \end{subfigure}&
        \begin{subfigure}{0.3\textwidth}
        \centering
        \caption{OpenAI3（$d=3072$）}
            \centering
            \includegraphics[width=\textwidth]{experiments/nearest_neighbor/recall-3072.pdf}
    \end{subfigure}
    \end{tabular}
    \caption{不同数据集与不同嵌入维度下的召回率对比。}
\end{figure}



\bibliography{paper}
\bibliographystyle{icml2025}

\newpage
\appendix
\onecolumn

\section{中英术语对照}
\label{sec:term-glossary}

为便于后续维护与复用，本文给出核心术语中英对照。后续增补内容建议优先沿用本表中的中文译法，以保证全文术语一致性。

\begin{table}[h]
\centering
\small
\setlength{\tabcolsep}{10pt}
\begin{tabular}{ll}
\toprule
\textbf{English Term} & \textbf{中文译法} \\
\midrule
Vector Quantization (VQ) & 向量量化 \\
Product Quantization (PQ) & 产品量化 \\
Distortion Rate & 失真率 \\
Mean-Squared Error (MSE) & 均方误差 \\
Inner Product Distortion / Error & 内积失真 / 内积误差 \\
Unbiased Estimator & 无偏估计器 \\
Bit-width & 位宽 \\
Codebook & 码本 \\
Centroid & 质心 \\
Residual Vector & 残差向量 \\
Random Rotation & 随机旋转 \\
Shannon Lower Bound (SLB) & Shannon 下界 \\
Yao's Minimax Principle & Yao 极小极大原理 \\
Near Neighbor (NN) Search & 近邻搜索 \\
Needle-In-A-Haystack & 大海捞针测试 \\
KV Cache Quantization & KV Cache 量化 \\
\bottomrule
\end{tabular}
\caption{论文核心术语中英对照表。}
\end{table}




\end{document}